\section{Část III: Specializace - Základy umělé inteligence}

\subsection{Řešení úloh prohledáváním}
\priorita{vysoká}{zelená (2 body)}

\begin{tema}
\textbf{Řešení úloh prohledáváním}
\begin{itemize}[leftmargin=1.2em,itemsep=0pt,topsep=1pt]
  \item formulace úlohy
  \item stromové vs. grafové prohledávání
  \item neinformované prohledávání (DFS, BFS, uniform-cost search)
  \item informované prohledávání (algoritmus A*, přípustné a konzistentní heuristiky)
\end{itemize}
\end{tema}

\ucivo
\begin{intuice}
Prohledávání je vhodné, když stav umíme brát jako uzel grafu a akci jako hranu. Odpověď u zkoušky má vždy říct: jak problém formuluji, jak držím frontier, podle čeho vybírám další uzel a kdy je řešení optimální.
\end{intuice}

\begin{center}
\begin{tikzpicture}[>=Stealth, node distance=9mm and 15mm,
  every node/.style={circle,draw,minimum size=7mm,font=\scriptsize,inner sep=0pt}]
  \node (s) {S};
  \node (a) [below left=of s] {A};
  \node (b) [below=of s] {B};
  \node (c) [below right=of s] {C};
  \node (d) [below left=of a] {D};
  \node (e) [below right=of a] {E};
  \node (g) [below=of c,fill=intugreen!15] {G};
  \draw[->] (s) -- node[above left,draw=none,font=\scriptsize]{$1$} (a);
  \draw[->] (s) -- node[right,draw=none,font=\scriptsize]{$4$} (b);
  \draw[->] (s) -- node[above right,draw=none,font=\scriptsize]{$2$} (c);
  \draw[->] (a) -- (d);
  \draw[->] (a) -- (e);
  \draw[->] (c) -- node[right,draw=none,font=\scriptsize]{$2$} (g);
  \node[draw=none,align=left,font=\scriptsize,anchor=west] at (3.0,-0.5)
    {BFS: S,A,B,C,D,E,G\\DFS: S,A,D,E,B,C,G\\UCS: podle $g(n)$\\A*: podle $g(n)+h(n)$};
\end{tikzpicture}
\end{center}

\begin{defbox}{formulace úlohy}
Úloha je dána počátečním stavem, množinou akcí nebo přechodovým modelem, testem cíle a cenou cesty. Stromové prohledávání si pamatuje cesty ve stromu možností, grafové prohledávání navíc uzavírá již navštívené stavy, aby se netočilo v cyklech.
\end{defbox}

\begin{itemize}[leftmargin=1.2em,itemsep=1pt,topsep=1pt]
  \item \textbf{BFS}: fronta, bere nejmělčí uzel. Je úplné pro konečné větvení a optimální při stejných cenách kroků.
  \item \textbf{DFS}: zásobník nebo rekurze, jde do hloubky. Má malou paměť, ale není obecně optimální a ve stromové verzi se může zacyklit.
  \item \textbf{Uniform-cost search}: prioritní fronta podle dosavadní ceny $g(n)$. Je Dijkstra nad stavovým prostorem, optimální pro nezáporné ceny (v nekonečném prostoru typicky pro kladnou dolní mez ceny kroku).
  \item \textbf{A$^*$}: vybírá minimum $f(n)=g(n)+h(n)$. Přípustná heuristika nikdy nepřeceňuje cenu do cíle. Konzistentní heuristika splňuje $h(n)\le c(n,a,n')+h(n')$ a dává optimalitu i v grafovém prohledávání bez znovuotevírání uzlů.
\end{itemize}
\klic{BFS = hloubka, DFS = paměť, UCS = skutečná cena, A$^*$ = skutečná cena plus odhad.}

\begin{odpoved}
\setlength{\parskip}{2pt}
\noindent\bod{1}\ \textbf{Minimum:} Prohledávací úloha je dána počátečním stavem, akcemi nebo přechodovým modelem, testem cíle a cenou cesty.
\emph{Přípustná heuristika} splňuje $0\le h(n)\le h^*(n)$, tedy nikdy nepřeceňuje skutečnou optimální cenu do cíle.\par
\noindent\bod{2}\ \textbf{+ k tomu:} \emph{Konzistentní heuristika} splňuje $h(n)\le c(n,a,n')+h(n')$ pro každý přechod a $h(G)=0$ pro cíl.
Tree search rozbaluje cesty ve stromu možností, graph search si pamatuje uzavřené stavy. BFS je optimální pro jednotkové ceny, UCS pro nezáporné ceny, A$^*$ pro přípustnou heuristiku ve stromu a konzistentní heuristiku v grafu.\par
\noindent\bod{3}\ \textbf{Bonus:} složitosti BFS $O(b^{d+1})$ čas i paměť, DFS čas $O(b^m)$ a paměť $O(bm)$ nebo $O(m)$ při backtrackingu. Dominance heuristiky: větší přípustná heuristika typicky expanduje méně uzlů.
Agent může být reflexní, modelový nebo cílový; atomické stavy se hodí pro prohledávání,
faktorizované pro CSP a plánování, strukturované pro predikátovou logiku. Greedy best-first
search používá jen $h(n)$ a není obecně optimální ani úplný.
\end{odpoved}

\begin{recall}
\begin{itemize}[leftmargin=1.2em,itemsep=0pt,topsep=1pt]
  \item Jakých pět částí má formulace problému prohledáváním?
  \item Kdy je BFS optimální a kdy nestačí?
  \item Jaký je rozdíl mezi přípustnou a konzistentní heuristikou?
  \item Co se stane, když DFS pustíš jako tree search na graf s cyklem?
\end{itemize}
\end{recall}
\past{Neříkej, že A$^*$ je vždy optimální. Potřebuje správnou heuristiku a u grafového prohledávání se typicky požaduje konzistence.}

\subsection{Splňování omezujících podmínek (CSP)}
\priorita{vysoká}{zelená (2 body)}

\begin{tema}
\textbf{Splňování omezujících podmínek}
\begin{itemize}[leftmargin=1.2em,itemsep=0pt,topsep=1pt]
  \item problém splňování podmínek
  \item hranová konzistence (algoritmus AC-3)
  \item algoritmus hledání řešení (MAC)
\end{itemize}
\end{tema}

\ucivo
\begin{intuice}
CSP neřeší cestu akcí, ale tabulku hodnot: každé proměnné chceme dát hodnotu tak, aby všechny podmínky seděly. Velký zisk je škrtat nemožné hodnoty dřív, než rozvětvíme backtracking.
\end{intuice}

\begin{center}
\begin{tikzpicture}[>=Stealth, node distance=16mm,
  var/.style={circle,draw,minimum size=8mm,font=\small},
  dom/.style={draw=none,font=\scriptsize,align=center}]
  \node[var] (a) {$A$};
  \node[var] (b) [right=of a] {$B$};
  \node[var] (c) [below=of a] {$C$};
  \draw (a) -- node[above,draw=none,font=\scriptsize]{$A<B$} (b);
  \draw (a) -- node[left,draw=none,font=\scriptsize]{$A\ne C$} (c);
  \draw (b) -- node[right,draw=none,font=\scriptsize]{$B\ne C$} (c);
  \node[dom,anchor=west] at (3.3,0.1) {před AC-3:\\$D_A=\{3,4,5\}$\\$D_B=\{1,2,3,4\}$};
  \node[dom,anchor=west] at (3.3,-1.25) {po AC pro $A<B$:\\$D_A=\{3\}$\\$D_B=\{4\}$};
\end{tikzpicture}
\end{center}

\begin{defbox}{CSP a hranová konzistence}
CSP má proměnné $X_i$, domény $D_i$ a podmínky jako relace na proměnných. Řešení je úplné přiřazení hodnot, které splní všechny podmínky. Hrana $(X_i,X_j)$ je hranově konzistentní, když pro každou hodnotu $x\in D_i$ existuje hodnota $y\in D_j$, která splní binární podmínku mezi $X_i$ a $X_j$.
\end{defbox}

\textbf{AC-3}: vlož všechny orientované hrany do fronty. Vyber $(X_i,X_j)$ a smaž z $D_i$ hodnoty bez podpory v $D_j$. Pokud se $D_i$ změnila, vrať do fronty všechny hrany $(X_k,X_i)$ kromě právě opačné. Prázdná doména znamená neúspěch. Jinak dostaneš lokálně vyčištěný CSP, ne nutně hotové řešení.

\textbf{MAC}: backtracking nad proměnnými, po každém přiřazení se znovu udržuje hranová konzistence. Tím se často odhalí spor dřív než hluboko ve stromu.

\begin{odpoved}
\setlength{\parskip}{2pt}
\noindent\bod{1}\ \textbf{Minimum:} \emph{CSP} je trojice proměnných $X_i$, domén $D_i$ a podmínek jako relací nad proměnnými; výstupem je úplné přiřazení hodnot z domén splňující všechny podmínky.
AC-3 škrtá hodnoty bez podpory u sousedních proměnných.\par
\noindent\bod{2}\ \textbf{+ k tomu:} Hrana $(X_i,X_j)$ je hranově konzistentní, když pro každé $x\in D_i$ existuje $y\in D_j$ splňující binární podmínku mezi $X_i$ a $X_j$.
MAC je backtracking plus opakované spuštění AC po každém přiřazení. Přidej mini příklad $A<B$.\par
\noindent\bod{3}\ \textbf{Bonus:} AC-3 pro binární podmínky má typicky složitost $O(cd^3)$, kde $c$ je počet podmínek a $d$ velikost domény. Hranová konzistence je jen lokální, proto sama negarantuje řešení.
Forward checking po přiřazení čistí jen domény dosud nepřiřazených sousedů, AC navíc
udržuje kompatibilitu mezi nepřiřazenými proměnnými. Obecnou $n$-ární podmínku lze převést
na binární za cenu pomocných proměnných.
\end{odpoved}

\begin{recall}
\begin{itemize}[leftmargin=1.2em,itemsep=0pt,topsep=1pt]
  \item Co přesně znamená, že hodnota má podporu v sousední doméně?
  \item Proč je hrana v CSP orientovaná, i když podmínka mezi proměnnými je binární?
  \item Jak se liší forward checking, AC-3 a MAC?
  \item Proč může být CSP hranově konzistentní a přesto nemít řešení?
\end{itemize}
\end{recall}
\past{Nezaměň hranovou konzistenci s globální konzistencí. AC-3 škrtá zjevně nemožné hodnoty, ale nemusí najít celé přiřazení.}

\subsection{Logické uvažování}
\priorita{vysoká}{zelená (2 body)}

\begin{tema}
\textbf{Logické uvažování}
\begin{itemize}[leftmargin=1.2em,itemsep=0pt,topsep=1pt]
  \item základy výrokové logiky (konjunktivní a disjunktivní normální forma)
  \item algoritmus DPLL
  \item dopředné a zpětné řetězení (Hornovské klauzule)
  \item rezoluce
\end{itemize}
\end{tema}

\ucivo
\begin{intuice}
Logické uvažování převádí znalosti na formule a otázku na test důsledku. Prakticky chceš umět dvě cesty: SAT přes DPLL a odvozování přes Hornovy klauzule nebo rezoluci.
\end{intuice}

\begin{defbox}{normální formy a klauzule}
CNF je konjunkce klauzulí, kde klauzule je disjunkce literálů. DNF je disjunkce konjunkcí literálů. Hornova klauzule má nejvýše jeden pozitivní literál, například $\neg p\lor\neg q\lor r$, což odpovídá pravidlu $p\land q\Rightarrow r$.
\end{defbox}

\textbf{DPLL} rozhoduje splnitelnost CNF. Opakuj jednotkovou propagaci, nastav čisté literály tak, aby splnily své klauzule, zkontroluj konec: žádné klauzule znamenají SAT, prázdná klauzule znamená UNSAT. Když není rozhodnuto, vyber proměnnou a větvi na pravda/nepravda.

\textbf{Řetězení}: dopředné řetězení jde od faktů k novým faktům, je data-driven. Zpětné řetězení jde od cíle k podcílům, je goal-driven. Obojí je levné a přirozené pro Hornovy klauzule.

\textbf{Rezoluce}: pro dotaz $\alpha$ testuj nesplnitelnost $KB\land\neg\alpha$ v CNF. Z klauzulí $(A\lor p)$ a $(B\lor\neg p)$ odvoď $(A\lor B)$. Když vznikne prázdná klauzule, $KB\models\alpha$.

\begin{odpoved}
\setlength{\parskip}{2pt}
\noindent\bod{1}\ \textbf{Minimum:} vysvětli CNF, DNF, klauzuli, literál, Hornovu klauzuli a že důsledek lze testovat přes $KB\land\neg\alpha$.\par
\noindent\bod{2}\ \textbf{+ k tomu:} popiš kroky DPLL: jednotková propagace, čistý literál, koncové stavy, větvení. Uveď rozdíl mezi dopředným a zpětným řetězením.\par
\noindent\bod{3}\ \textbf{Bonus:} napiš rezoluční pravidlo a vysvětli rezoluční zamítnutí. DPLL je v nejhorším exponenciální, ale jednotková propagace drasticky zmenšuje praxi.
Dopředné řetězení lze implementovat počítadly nesplněných předpokladů pravidel; zpětné
řetězení může být výrazně levnější, když dotaz míří jen na malou část báze znalostí.
\end{odpoved}

\begin{recall}
\begin{itemize}[leftmargin=1.2em,itemsep=0pt,topsep=1pt]
  \item Co udělá DPLL s jednotkovou klauzulí?
  \item Jak přepíšeš $\neg p\lor\neg q\lor r$ jako pravidlo?
  \item Proč se při důkazu dotazu přidává $\neg\alpha$?
  \item Kdy je výhodnější zpětné řetězení než dopředné?
\end{itemize}
\end{recall}
\past{Nepleť DPLL a rezoluci: DPLL hledá ohodnocení CNF, rezoluce odvozuje nové klauzule a hledá prázdnou klauzuli.}

\subsection{Pravděpodobnostní uvažování}
\priorita{vysoká}{zelená (2 body)}

\begin{tema}
\textbf{Pravděpodobnostní uvažování}
\begin{itemize}[leftmargin=1.2em,itemsep=0pt,topsep=1pt]
  \item základy teorie pravděpodobnosti (úplná sdružená distribuce, nezávislost, Bayesovo pravidlo)
  \item pravděpodobnostní uvažování (vysčítání, normalizace)
  \item Bayesovské sítě
  \item konstrukce a vztah k úplné sdružené distribuci
  \item exaktní odvozování (enumerace, eliminace proměnných)
  \item aproximační odvozování (Monte Carlo, zamítání, vážení věrohodností)
\end{itemize}
\end{tema}

\ucivo
\begin{intuice}
Pravděpodobnost je způsob, jak uvažovat při nejistotě. Úplná sdružená distribuce umí odpovědět na všechno, ale je obrovská. Bayesovská síť ji komprimuje pomocí podmíněných nezávislostí.
\end{intuice}

\begin{center}
\begin{tikzpicture}[>=Stealth, node distance=11mm and 16mm,
  every node/.style={draw,rounded corners=1pt,minimum width=12mm,minimum height=6mm,font=\scriptsize}]
  \node (b) {Vloupání};
  \node (e) [right=of b] {Zemětř.};
  \node (a) [below right=of b] {Alarm};
  \node (j) [below left=of a] {Jan volá};
  \node (m) [below right=of a] {Marie volá};
  \draw[->] (b) -- (a);
  \draw[->] (e) -- (a);
  \draw[->] (a) -- (j);
  \draw[->] (a) -- (m);
  \node[draw=none,align=left,anchor=west,font=\scriptsize] at (4.1,-0.8)
    {$P(B,E,A,J,M)=$\\$P(B)P(E)P(A\mid B,E)$\\$\cdot P(J\mid A)P(M\mid A)$};
\end{tikzpicture}
\end{center}

\begin{defbox}{Bayesovská síť}
Bayesovská síť je orientovaný acyklický graf náhodných proměnných a tabulky $P(X_i\mid Parents(X_i))$. Společná distribuce se faktorizuje jako
\[
P(x_1,\dots,x_n)=\prod_i P(x_i\mid pa_i),
\]
kde $pa_i$ jsou hodnoty rodičů proměnné $X_i$.
\end{defbox}

\textbf{Základní počty}: z úplné sdružené distribuce získáš marginál vysčítáním skrytých proměnných, například $P(X)=\sum_y P(X,y)$. Podmíněnou pravděpodobnost získáš normalizací: $P(X\mid e)=\alpha P(X,e)$. Bayesovo pravidlo je $P(H\mid e)=\alpha P(e\mid H)P(H)$.

\textbf{Inference v síti}: enumerace přímo sčítá přes skryté proměnné. Eliminace proměnných průběžně tvoří faktory a sčítá jen jednou, takže šetří opakovanou práci. Aproximace generuje vzorky: Monte Carlo obecně, zamítání zahodí vzorky neslučitelné s evidencí, vážení věrohodností vzorky nezahazuje, ale váží je pravděpodobností evidence.

\begin{odpoved}
\setlength{\parskip}{2pt}
\noindent\bod{1}\ \textbf{Minimum:} \emph{Bayesovská síť} je DAG náhodných proměnných s tabulkami $P(X_i\mid Parents(X_i))$; společnou distribuci faktorizuje jako $\prod_i P(x_i\mid pa_i)$.
Podmíněná nezávislost znamená $P(X\mid Y,Z)=P(X\mid Z)$.\par
\noindent\bod{2}\ \textbf{+ k tomu:} Marginál získáš vysčítáním skrytých proměnných, podmíněnou pravděpodobnost normalizací $P(X\mid e)=\alpha P(X,e)$.
Popiš enumeraci, eliminaci proměnných a rozdíl mezi zamítáním a vážením věrohodností.\par
\noindent\bod{3}\ \textbf{Bonus:} vysvětli, že konstrukce sítě závisí na pořadí proměnných a nejlepší je obvykle směr od příčin k důsledkům. Eliminace závisí na pořadí eliminovaných proměnných.
Chain rule rozkládá sdruženou distribuci na součin podmíněných pravděpodobností.
Naivní Bayes předpokládá podmíněnou nezávislost důsledků při známé příčině.
Likelihood weighting fixuje evidenci a váží vzorky pravděpodobností evidence.
\end{odpoved}

\begin{recall}
\begin{itemize}[leftmargin=1.2em,itemsep=0pt,topsep=1pt]
  \item Jak dostaneš $P(X\mid e)$ z $P(X,e)$?
  \item Co přesně faktorizuje Bayesovská síť?
  \item Proč je eliminace proměnných úspornější než čistá enumerace?
  \item Jaký je rozdíl mezi zamítáním a vážením věrohodností?
\end{itemize}
\end{recall}
\past{Neříkej, že chybějící hrana znamená obyčejnou nezávislost. V Bayesovské síti jde hlavně o podmíněné nezávislosti vzhledem k rodičům a evidenci.}

\subsection{Reprezentace znalostí}
\priorita{vysoká}{zelená (2 body)}

\begin{tema}
\textbf{Reprezentace znalostí}
\begin{itemize}[leftmargin=1.2em,itemsep=0pt,topsep=1pt]
  \item situační kalkulus, problém rámce
  \item Markovské modely
  \item filtrace, predikce, vyhlazování, nejpravděpodobnější průchod
  \item skryté Markovské modely (HMM) vs. dynamické Bayesovské sítě
\end{itemize}
\end{tema}

\ucivo
\begin{intuice}
Reprezentace znalostí řeší, jak zapsat svět tak, aby nad ním šlo odvozovat. Situační kalkulus je logický zápis změn po akcích, Markovské modely jsou pravděpodobnostní zápis vývoje stavu v čase.
\end{intuice}

\begin{center}
\begin{tikzpicture}[>=Stealth, node distance=12mm and 13mm,
  hidden/.style={circle,draw,minimum size=7mm,font=\scriptsize},
  obs/.style={rectangle,draw,minimum size=6mm,font=\scriptsize}]
  \node[hidden] (x1) {$X_1$};
  \node[hidden] (x2) [right=of x1] {$X_2$};
  \node[hidden] (x3) [right=of x2] {$X_3$};
  \node[obs] (e1) [below=of x1] {$E_1$};
  \node[obs] (e2) [below=of x2] {$E_2$};
  \node[obs] (e3) [below=of x3] {$E_3$};
  \draw[->] (x1) -- (x2);
  \draw[->] (x2) -- (x3);
  \draw[->] (x1) -- (e1);
  \draw[->] (x2) -- (e2);
  \draw[->] (x3) -- (e3);
  \node[draw=none,anchor=west,align=left,font=\scriptsize] at (4.0,-0.7)
    {HMM:\\skrytý stav $X_t$\\pozorování $E_t$};
\end{tikzpicture}
\end{center}

\begin{defbox}{situační kalkulus a HMM}
Situační kalkulus popisuje situace vzniklé posloupností akcí, predikáty o situacích a akci $Result(s,a)$. Problém rámce je otázka, jak stručně zapsat, co se akcí nemění. HMM má skrytý stav $X_t$, pozorování $E_t$, přechod $P(X_t\mid X_{t-1})$ a senzorový model $P(E_t\mid X_t)$.
\end{defbox}

\textbf{Markovské úlohy}: filtrace počítá $P(X_t\mid e_{1:t})$, tedy aktuální stav po dosavadních pozorováních. Predikce počítá budoucí stav $P(X_{t+k}\mid e_{1:t})$. Vyhlazování zpřesňuje minulý stav $P(X_k\mid e_{1:t})$ pro $k<t$. Nejpravděpodobnější průchod hledá nejpravděpodobnější posloupnost skrytých stavů, typicky Viterbiho algoritmem.

\textbf{HMM vs. dynamická Bayesovská síť}: HMM má jednu skrytou stavovou proměnnou na čas. Dynamická Bayesovská síť rozkládá stav na více proměnných a dovolí bohatší závislosti mezi nimi.

\begin{odpoved}
\setlength{\parskip}{2pt}
\noindent\bod{1}\ \textbf{Minimum:} \emph{HMM} je model se skrytými stavy $X_t$, pozorováními $E_t$, počáteční distribucí, přechody $P(X_t\mid X_{t-1})$ a senzorem $P(E_t\mid X_t)$.
Situační kalkulus popisuje situace vzniklé posloupností akcí a funkci $Result(s,a)$.\par
\noindent\bod{2}\ \textbf{+ k tomu:} Markovův předpoklad: $X_t$ závisí na minulosti jen přes $X_{t-1}$; pozorování závisí jen na aktuálním stavu.
Filtrace počítá $P(X_t\mid e_{1:t})$, predikce $P(X_{t+k}\mid e_{1:t})$, vyhlazování $P(X_k\mid e_{1:t})$ pro $k<t$ a Viterbi nejpravděpodobnější průchod.\par
\noindent\bod{3}\ \textbf{Bonus:} uveď, že HMM je speciální případ dynamické Bayesovské sítě; DBN dovolí faktorizovat stav na více proměnných a tím zmenšit tabulky.
Markovův předpoklad říká, že další stav závisí jen na předchozím; stacionarita znamená
stejné přechodové tabulky v čase. Filtrace používá forward algoritmus, vyhlazování
forward-backward a nejpravděpodobnější průchod Viterbiho algoritmus.
\end{odpoved}

\begin{recall}
\begin{itemize}[leftmargin=1.2em,itemsep=0pt,topsep=1pt]
  \item Co je problém rámce v jedné větě?
  \item Jaký dotaz odpovídá filtraci a jaký vyhlazování?
  \item Co hledá Viterbiho algoritmus?
  \item Proč je DBN obecnější než HMM?
\end{itemize}
\end{recall}
\past{Nezaměň filtraci s vyhlazováním: filtrace odhaduje současnost z minulých pozorování, vyhlazování zpětně zpřesňuje minulost i pomocí pozdější evidence.}

\subsection{Automatické plánování}
\priorita{vysoká}{zelená (2 body)}

\begin{tema}
\textbf{Automatické plánování}
\begin{itemize}[leftmargin=1.2em,itemsep=0pt,topsep=1pt]
  \item formulace plánovacího problému (definice operátoru)
  \item dopředné a zpětné plánování
\end{itemize}
\end{tema}

\ucivo
\begin{intuice}
Plánování je prohledávání nad faktorizovaným stavem. Místo černé skříňky používáme fakta, předpoklady akcí a efekty akcí, takže lze rozumněji vybírat relevantní kroky.
\end{intuice}

\begin{defbox}{plánovací problém a operátor}
Klasický plánovací problém má počáteční stav jako množinu faktů, cíl jako podmínku na fakta a operátory. Operátor má preconditions, tedy kdy ho lze použít, a effects, tedy co přidá nebo smaže ve stavu. Plán je posloupnost operátorů, která z počátečního stavu vyrobí stav splňující cíl.
\end{defbox}

Mini příklad: operátor \emph{Přesuň(x,A,B)} má předpoklady $Na(x,A)$ a $Volne(B)$, efekty přidat $Na(x,B)$ a smazat $Na(x,A)$. Dopředné plánování začíná v počátečním stavu a zkouší použitelné akce. Zpětné plánování začíná cílem a hledá akce, které by cíl mohly splnit, přičemž jejich předpoklady se stanou novými podcíli.

\begin{itemize}[leftmargin=1.2em,itemsep=1pt,topsep=1pt]
  \item \textbf{Dopředné}: jednoduché simulování akcí, ale může zkoušet mnoho irelevantních kroků.
  \item \textbf{Zpětné}: soustředí se na cíl, ale musí řešit, zda vybrané akce jsou vzájemně slučitelné.
  \item \textbf{Kdy použít}: dopředné je přirozené, když je málo použitelných akcí; zpětné, když cíl silně omezuje relevantní akce.
\end{itemize}

\begin{odpoved}
\setlength{\parskip}{2pt}
\noindent\bod{1}\ \textbf{Minimum:} plánovací problém = počáteční stav, cílová podmínka, operátory. Operátor = předpoklady a efekty. Plán = posloupnost akcí.\par
\noindent\bod{2}\ \textbf{+ k tomu:} vysvětli dopředné a zpětné plánování na mini příkladu s přesunem objektu. Řekni, proč faktorizovaný stav umožní pracovat s relevancí akcí.\par
\noindent\bod{3}\ \textbf{Bonus:} uveď rozdíl proti obyčejnému prohledávání: stav není atomický, akce mají logickou strukturu a plánovač může využívat předpoklady a efekty místo slepé expanze.
Plánování používá předpoklad uzavřeného světa; fluenty se akcemi mění, rigidní atomy ne.
Akce je plně instanciovaný operátor a zpětné plánování používá jen akce relevantní pro cíl.
Situační kalkulus řeší změny přes predikát proveditelnosti $Poss(a,s)$ a successor-state axiomy.
\end{odpoved}

\begin{recall}
\begin{itemize}[leftmargin=1.2em,itemsep=0pt,topsep=1pt]
  \item Jaké části má operátor?
  \item Co je stav v klasickém plánování?
  \item Jak se liší dopředné a zpětné plánování?
  \item Proč je plánování vhodné pro faktorizované stavy?
\end{itemize}
\end{recall}
\past{Neodpovídej jen obecně "hledání cesty". U plánování musí zaznít preconditions, effects a cíl jako podmínka na fakta.}

\subsection{Markovské rozhodovací procesy (MDP)}
\priorita{vysoká}{zelená (2 body)}

\begin{tema}
\textbf{Markovské rozhodovací procesy (MDP)}
\begin{itemize}[leftmargin=1.2em,itemsep=0pt,topsep=1pt]
  \item formulace problému (výpočet užitku, strategie)
  \item Bellmanova rovnice
  \item iterace hodnot, iterace strategií
  \item POMDP (základní definice)
\end{itemize}
\end{tema}

\ucivo
\begin{intuice}
MDP je plánování, kde akce nejsou jisté. Agent volí akci podle stavu, dostane odměnu a pravděpodobnostně přejde do dalšího stavu. Hledáme strategii s nejvyšším očekávaným diskontovaným užitkem.
\end{intuice}

\begin{center}
\begin{tikzpicture}[>=Stealth, node distance=22mm,
  st/.style={circle,draw,minimum size=9mm,font=\small}]
  \node[st] (s0) {$S_0$};
  \node[st] (s1) [right=of s0] {$S_1$};
  \node[st,fill=intugreen!15] (g) [right=of s1] {$G$};
  \draw[->] (s0) to[bend left=15] node[above,draw=none,font=\scriptsize]{$a:0.8, r=1$} (s1);
  \draw[->] (s0) to[loop above] node[above,draw=none,font=\scriptsize]{$0.2$} (s0);
  \draw[->] (s1) -- node[above,draw=none,font=\scriptsize]{$a:1, r=5$} (g);
  \draw[->] (s1) to[loop below] node[below,draw=none,font=\scriptsize]{$b:1, r=0$} (s1);
\end{tikzpicture}
\end{center}

\begin{defbox}{MDP}
MDP je čtveřice stavů, akcí, přechodových pravděpodobností $P(s'\mid s,a)$ a odměn $R(s,a,s')$, často s diskontem $\gamma$. Strategie $\pi(s)$ říká, jakou akci zvolit ve stavu $s$. Užitečnost stavu je očekávaný součet budoucích diskontovaných odměn.
\end{defbox}

\textbf{Bellmanova rovnice optimality}:
\[
U(s)=\max_a \sum_{s'} P(s'\mid s,a)\bigl(R(s,a,s')+\gamma U(s')\bigr).
\]
\textbf{Iterace hodnot}: opakovaně aktualizuj $U(s)$ Bellmanovou rovnicí a nakonec z ní vyčti strategii. \textbf{Iterace strategií}: střídá vyhodnocení současné strategie a její zlepšení výběrem lepší akce. \textbf{POMDP}: stav není přímo pozorovaný, agent udržuje belief, tedy distribuci přes možné stavy.

\begin{odpoved}
\setlength{\parskip}{2pt}
\noindent\bod{1}\ \textbf{Minimum:} \emph{MDP} je pětice stavů $S$, akcí $A$, přechodů $P(s'\mid s,a)$, odměn $R(s,a,s')$ a diskontu $\gamma$; strategie $\pi(s)$ vybírá akci ve stavu.
Cíl je maximalizovat očekávaný diskontovaný užitek.\par
\noindent\bod{2}\ \textbf{+ k tomu:} Bellmanova rovnice optimality:
$U(s)=\max_a\sum_{s'}P(s'\mid s,a)(R(s,a,s')+\gamma U(s'))$.
Iterace hodnot opakovaně aktualizuje $U$, iterace strategií střídá vyhodnocení a zlepšení strategie.\par
\noindent\bod{3}\ \textbf{Bonus:} POMDP = MDP s částečnou pozorovatelností, rozhoduje se nad belief stavem. Iterace strategií bývá dražší v kroku vyhodnocení, ale často potřebuje méně iterací.
Diskontovaný užitek trajektorie je konečný pro $\gamma<1$ a omezené odměny. Optimální
strategie závisí na aktuálním stavu, ne na počátečním. Policy evaluation řeší Bellmanovy
rovnice bez maxima, policy improvement vybírá lepší akce.
\end{odpoved}

\begin{recall}
\begin{itemize}[leftmargin=1.2em,itemsep=0pt,topsep=1pt]
  \item Co je strategie v MDP?
  \item Kde v Bellmanově rovnici vystupuje nejistota?
  \item Jak se liší iterace hodnot a iterace strategií?
  \item Co je belief stav u POMDP?
\end{itemize}
\end{recall}
\past{Nepleť MDP s HMM: v MDP agent volí akce a optimalizuje odměny, v HMM hlavně odhaduje skrytý stav z pozorování.}

\subsection{Hry a teorie her}
\priorita{vysoká}{zelená (2 body)}

\begin{tema}
\textbf{Hry a teorie her}
\begin{itemize}[leftmargin=1.2em,itemsep=0pt,topsep=1pt]
  \item algoritmy Minimax a alfa-beta prořezávání
  \item základy teorie her (vězňovo dilema, Nashovo ekvilibrium)
  \item mechanism design (typy aukcí)
\end{itemize}
\end{tema}

\ucivo
\begin{intuice}
Ve hrách neoptimalizuješ proti přírodě, ale proti jinému agentovi. Minimax předpokládá, že soupeř hraje proti nám co nejlépe. Alfa-beta jen šetří práci, výsledek minimaxu nemění.
\end{intuice}

\begin{center}
\begin{tikzpicture}[>=Stealth, level distance=10mm,
  level 1/.style={sibling distance=34mm},
  level 2/.style={sibling distance=15mm},
  every node/.style={font=\scriptsize}]
  \node[circle,draw] {MAX}
    child {node[circle,draw] {MIN}
      child {node[draw] {$3$}}
      child {node[draw] {$5$}}
    }
    child {node[circle,draw] {MIN}
      child {node[draw] {$2$}}
      child {node[draw,gray] {$9$} edge from parent[gray,dashed]}
      edge from parent node[above,draw=none]{$\alpha=3$}
    };
  \node[draw=none,align=left,anchor=west] at (2.6,-1.0)
    {vpravo: po hodnotě 2\\MIN už nemůže zlepšit MAX\\větev 9 se řeže};
\end{tikzpicture}
\end{center}

\begin{defbox}{minimax, Nashovo ekvilibrium}
Minimax počítá hodnotu hry s nulovým součtem: v uzlu MAX bere maximum z potomků, v uzlu MIN minimum. Nashovo ekvilibrium je profil strategií, kde se žádnému hráči nevyplatí jednostranně změnit strategii.
\end{defbox}

\textbf{Alfa-beta}: drží $\alpha$ jako nejlepší zatím zajištěnou hodnotu pro MAX a $\beta$ jako nejlepší zatím zajištěnou hodnotu pro MIN. Když je jasné, že větev nemůže ovlivnit rozhodnutí předka, neprohledává se.

\textbf{Vězňovo dilema}: každý má individuální motivaci zradit, i když oboustranná spolupráce je společně lepší. Ukazuje rozdíl mezi individuální racionalitou a kolektivním výsledkem.

\textbf{Aukce}: anglická je vzestupná, holandská sestupná, první cena znamená platbu vlastní nabídky, druhá cena typicky motivuje nabídnout skutečnou hodnotu.

\begin{odpoved}
\setlength{\parskip}{2pt}
\noindent\bod{1}\ \textbf{Minimum:} minimax = MAX maximalizuje, MIN minimalizuje. Alfa-beta prořezává větve bez změny výsledku. Nash = žádná jednostranně výhodná odchylka.\par
\noindent\bod{2}\ \textbf{+ k tomu:} ukaž výpočet malého stromu a jeden alfa-beta řez. Vysvětli vězňovo dilema a vyjmenuj základní typy aukcí.\par
\noindent\bod{3}\ \textbf{Bonus:} dobré pořadí tahů dává alfa-beta velkou úsporu, v ideálním případě zhruba na efektivní hloubku poloviny stromu. U aukce druhé ceny zmiň pravdivé nabízení jako dominantní strategii v jednoduchém modelu.
V hrách v normálním tvaru rozliš čisté a smíšené strategie, dominantní strategii,
Pareto optimalitu a Nashovo ekvilibrium. Vickreyho aukce je obálková aukce druhé ceny;
VCG zobecňuje pravdivé mechanismy pro víceparametrické aukce.
\end{odpoved}

\begin{recall}
\begin{itemize}[leftmargin=1.2em,itemsep=0pt,topsep=1pt]
  \item Jak se počítá hodnota v MAX a MIN uzlu?
  \item Proč alfa-beta nemění výsledek minimaxu?
  \item Co znamená Nashovo ekvilibrium?
  \item Jaký je rozdíl mezi aukcí první a druhé ceny?
\end{itemize}
\end{recall}
\past{Neříkej, že alfa-beta je jiná strategie hraní. Je to optimalizace výpočtu stejného minimaxového rozhodnutí.}

\subsection{Strojové učení (přehledově)}
\priorita{vysoká}{zelená (2 body)}

\begin{tema}
\textbf{Strojové učení}
\begin{itemize}[leftmargin=1.2em,itemsep=0pt,topsep=1pt]
  \item základní druhy učení (s učitelem, bez učitele, zpětnovazební)
  \item rozhodovací stromy (definice, konstrukce)
  \item regrese, SVM (základní principy)
  \item Bayesovské učení, EM algoritmus
  \item zpětnovazební učení
  \item pasivní učení (definice, metody ADP a TD)
  \item aktivní učení (definice, explorace vs. exploitace, Q-učení, SARSA)
\end{itemize}
\end{tema}

\ucivo
\begin{intuice}
Tady se nečeká detail celé specializace Strojové učení, ale mapa pojmů v kontextu UI. Bezpečná odpověď vždy řekne typ úlohy, co se učí, z jakých dat a jak se pozná dobrý výsledek.
\end{intuice}

\begin{defbox}{základní druhy učení}
Učení s učitelem má dvojice vstup plus správný výstup a řeší klasifikaci nebo regresi. Učení bez učitele hledá strukturu bez cílových štítků. Zpětnovazební učení má agenta, akce, odměny a učí strategii z interakce s prostředím.
\end{defbox}

\textbf{Rozhodovací strom}: vnitřní uzel testuje atribut, hrana odpovídá výsledku testu a list dává predikci. Konstrukce vybírá test, který nejlépe rozdělí data, typicky podle informačního zisku nebo Gini nečistoty, a zastavuje nebo prořezává kvůli přeučení.

\textbf{Regrese a SVM}: regrese predikuje spojitou hodnotu, typicky minimalizuje chybu. SVM hledá oddělující nadrovinu s maximálním marginem; pro neseparabilní data používá měkký margin a pro nelineární hranice jádra.

\textbf{Bayesovské učení a EM}: Bayesovské učení kombinuje prior a věrohodnost na posterior $P(h\mid data)$. EM řeší skryté proměnné střídáním E-kroku, který dopočítá očekávání skrytých přiřazení, a M-kroku, který zlepší parametry.

\textbf{Zpětnovazební učení}: pasivní učení hodnotí pevně danou strategii. ADP odhaduje model přechodů a odměn a pak řeší MDP. TD aktualizuje hodnoty z rozdílu mezi po sobě jdoucími odhady. Aktivní učení současně volí akce: musí vyvažovat exploraci a exploitaci. Q-učení je off-policy aktualizace hodnot akcí, SARSA je on-policy aktualizace podle skutečně provedeného přechodu $(s,a,r,s',a')$.

\begin{odpoved}
\setlength{\parskip}{2pt}
\noindent\bod{1}\ \textbf{Minimum:} Učení s učitelem má dvojice vstup-výstup a učí predikci; učení bez učitele hledá strukturu bez štítků; zpětnovazební učení hledá strategii maximalizující očekávanou odměnu.
Rozhodovací strom testuje atributy ve vnitřních uzlech a predikuje v listech; SVM hledá nadrovinu s maximálním marginem.\par
\noindent\bod{2}\ \textbf{+ k tomu:} popiš konstrukci stromu, princip marginu u SVM, Bayesovské učení přes posterior, EM jako E/M kroky, pasivní ADP/TD a aktivní Q-učení/SARSA.\par
\noindent\bod{3}\ \textbf{Bonus:} zdůrazni rozdíl ADP vs. TD: ADP se učí model a řeší ho, TD se učí hodnoty přímo. Q-učení je off-policy s maximem přes další akce, SARSA je on-policy a používá další skutečně zvolenou akci.
Přímý odhad utility průměruje reward-to-go, ale ignoruje Bellmanovy vazby. TD update:
$U(s)\leftarrow U(s)+\alpha(R(s)+\gamma U(s')-U(s))$. Pro aktivní RL je exploration
nutná, protože hladový agent může zůstat u špatně naučené akce; $Q(s,a)$ dovolí volit akci
bez znalosti přechodového modelu.
\end{odpoved}

\begin{recall}
\begin{itemize}[leftmargin=1.2em,itemsep=0pt,topsep=1pt]
  \item Jak rozlišíš klasifikaci a regresi?
  \item Co je list a co vnitřní uzel v rozhodovacím stromu?
  \item Co se děje v E-kroku a M-kroku EM algoritmu?
  \item Jaký je rozdíl mezi Q-učením a SARSA?
\end{itemize}
\end{recall}
\past{Nezajížděj mimo oficiální přehled: žádné hluboké učení, NLP ani robotika. Tady stačí základní principy uvedených metod.}
